\chapter{Results}
\label{ch:resultsmain}

Every number below carries its strength of evidence, and \textbf{no result
before the pilot in \cref{sec:pilot-results} was validated on a physical
Kinova arm carrying a load}: most of this project is simulated against mock
hardware, using the delivered system's own kinematic description, solver
and safety code, plus recorded data from the physical master and four
cameras (\cref{ch:results}: full derivations).

\section{The command mapping and its safety layer}

Channel health, computed over \emph{distinct} sensor updates rather than
transmitted rows (roughly \SIrange{82}{88}{\percent} of consecutive rows
repeat, which had previously been misread as dead channels), found six of
fourteen channels alive (\cref{fig:channelhealth}).

\begin{figure}[htbp]
  \centering
  \includegraphics[width=0.92\linewidth]{channels.pdf}
  \caption[Master channel health, measured]{Every master channel; bar
  colour is the verdict. Above the five per cent line (upper panel), a
  channel is fresh but not tracing a trajectory: \emph{incoherent}.}
  \label{fig:channelhealth}
\end{figure}

Replaying the recorded incoherent traces with the operator's input held
constant shows those channels alone inject \SI{54.5}{\milli\metre} of
commanded motion on the left arm and \SI{0.014}{\milli\metre} on the right
-- motion nobody asked for. Regressing true reach on the surviving channels
shows the left arm retains no observable radial dimension at all
($R^2=\num{0.133}$) against the right's near-complete
$R^2=\num{0.986}$, naming the repair precisely as two left-arm channels.

Vertical and fore-aft command axes track correctly on both arms. Lateral
does not: two master motions meant \SI{90}{\degree} apart arrive
\SI{137.7}{\degree} apart, and since a rotation preserves angles, no
remapping turns a non-orthogonal measured triad orthogonal -- a geometric
limitation, not a calibration defect.

The safety layer handled fourteen of fourteen injected faults (channel
dropout, board disconnect/reconnect, emergency stop mid-motion,
wrong-object latching, stale detections, session loss, clock
non-monotonicity and more). Two needed several attempts, both the same
class of error: a mechanism keyed on the wrong observable. The dead-man did
not trip on ten seconds of total master silence, because the master
substitutes stale data at full rate; keying it on the data's source
timestamp fixed it. The same path once called a blocking driver shutdown
first, adding \SI{4}{\second} between deciding to stop and the arm
stopping; the fix parks the arm first, from local publishes only.

\section{Workspace and motion}

Over a 63-cell grid spanning the frontal working volume with both arms
verified at the home pose, \textbf{no cell was reachable by both arms}
simultaneously (\cref{fig:disjoint}), and the result survives freeing the
wrist yaw and adding solver retries.

\begin{figure}[htbp]
  \centering
  \begin{subfigure}[b]{0.48\linewidth}
    \centering\includegraphics[width=\linewidth]{workspace_disjoint.pdf}
    \caption{the measured dead band}
    \label{fig:disjoint-a}
  \end{subfigure}\hfill
  \begin{subfigure}[b]{0.48\linewidth}
    \centering\includegraphics[width=\linewidth]{mount_sweep.pdf}
    \caption{the costed fix}
    \label{fig:disjoint-b}
  \end{subfigure}
  \caption[The disjoint workspace, and the mount change that fixes it]{
  (\subref{fig:disjoint-a}) The two reachable sets do not meet. (\subref{fig:disjoint-b}) A mount-placement sweep: forward tilt and inboard
  translation help, splaying outward \emph{hurts}. Kinematic only (the
  wearer is not moved), so these are upper bounds pending a clearance check.}
  \label{fig:disjoint}
\end{figure}

Moving the mounts inboard, forward and down with extra forward tilt raises
shared cells to 22 of 63 within a buildable \SI{0.20}{\metre} separation.
That change was applied, raising home-pose clearance from
\SI{0.1610}{\metre} (a reading of the mount itself, since no joint moves
it) to \SI{0.2202}{\metre}, and moving pick-and-place targets from
\SI{0.530}{\metre} off centre to \SI{0.290}{\metre}. Neither arm crosses
the body centre line at any surface height or distance in \num{3360} swept
configurations; the constraint is the wearer's torso, not arm posture --
deleting the wearer's arms from the model entirely moves the boundary by at
most \SI{75}{\milli\metre}.

The per-joint step limiter that constituted the whole of teleoperation
motion generation, measured against a \SI{30}{\milli\metre} grasp capture
gate, is compared with its jerk-limited replacement in
\cref{tab:motiongen-main}: it left the hand off the solved path by nearly
twice the gate, at over twelve times the manufacturer's declared joint
speed, because it never read the limits file.

\begin{table}[htbp]
  \centering
  \small
  \caption[Motion generator comparison]{The same commanded slew, both
  arms, at the master's own \SI{50}{\hertz}. Off-path is the largest
  distance the hand leaves the solver's straight line.}
  \label{tab:motiongen-main}
  \begin{tabular}{@{}lrrr@{}}
    \toprule
    Generator & Off-path & Peak $\dot q$ of limit & Arrival \\
    \midrule
    Per-joint clamp, as shipped & \SI{51.3}{\milli\metre} & $12.53\times$ & 2 cycles \\
    \textbf{Jerk-limited, delivered} & \textbf{\SI{0.0}{\milli\metre}} & $1.00\times$ & 1 cycle \\
    \bottomrule
  \end{tabular}
\end{table}

Replaying recorded joint commands on the physical arm shows every joint
parks \SI{0.305}{\degree} short of its target, on the side it approached
from -- a terminal offset, confirmed by propagating the error through the
manipulator Jacobian ($r=\num{0.923}$) rather than a proportional-gain
model. A single scalar per arm, applied as overshoot in the direction of
travel, removes \SIrange{74}{76}{\percent} of it. Folded into the
positioning error budget with a re-derived pad offset and a
measured-surface centring correction, this takes the combined budget from
\SI{18.64}{} to \SI{12.91}{\milli\metre} RSS and \SI{33.39}{} to
\SI{19.99}{\milli\metre} worst case -- inside the \SI{30}{\milli\metre}
requirement, \emph{provided} the still-unmeasured camera-to-wrist transform
is small. The first sustained attempt to drive a physical arm from its own
wrist camera measured that transform directly and found it is not: four
poses of a stationary table disagreed by \SI{8.45}{\degree}, \SIrange{34}{59}{\milli\metre}
of position error over the working range (\cref{sec:hardware-session}).

\section{Perception: what the cameras actually detected}

Instrumented at thirty-three stages per camera and pointed at four physical
cameras for the first time (\cref{ch:vision}), the wrist cameras fit a
support plane at \SI{1.43}{\milli\metre} RMS with \SI{74}{\percent} inliers
at working distance (\cref{fig:vision-plane}) -- within a factor of three of
the same code against a renderer, and the first evidence the surface-fitting
stage transfers to hardware.

\begin{figure}[htbp]
  \centering
  \includegraphics[width=0.92\linewidth]{vision/wrist_ransac.png}
  \caption[Detection: the support plane, fitted and scored on a real
  camera]{The support plane fitted by random sample consensus on the
  left-hand wrist camera. Green is an inlier (on the fitted plane), red is
  not. The inlier fraction is the quantity that matters, not the residual
  alone: a plane can fit three points perfectly.}
  \label{fig:vision-plane}
\end{figure}

Objects are recovered as oriented three-dimensional boxes rather than an
axis-aligned bounding volume (\cref{fig:vision-boxes}), which is what makes
a turned object's true width -- the quantity the grasp planner needs --
available at all.

\begin{figure}[htbp]
  \centering
  \includegraphics[width=0.92\linewidth]{vision/wrist_boxes3d.png}
  \caption[Detection: objects as oriented 3D boxes, projected and
  plan-view]{Five objects recovered from the left-hand wrist camera, each
  posed as an oriented box and projected back into the picture (top) and
  shown in plan view (bottom).}
  \label{fig:vision-boxes}
\end{figure}

No grasp closed on any camera this session: the nearest object's narrowest
axis measured \SI{144}{\milli\metre} against an \SI{85}{\milli\metre} jaw, a
statement about the scene rather than the planner. The same session found,
within its first hour, that three files independently asserted the scene
camera's mounting orientation and none had measured it, silently inverting
every earlier scene figure -- invisible to every geometric metric, since a
consistent rotation preserves every distance, and found only by comparing
one camera's picture against another's.

In simulation, the full task set solves \num{13927} inverse-kinematics
calls with zero failures and zero clearance breaches, reproducing field for
field on re-run. Sweeping the shared-autonomy intent estimator against a
mis-declared table shows why this matters operationally: from declared
coordinates it becomes \emph{confidently wrong} \SI{50}{\percent} of the
time once the declaration is one object-pitch out; from measured
coordinates it stays at \SI{0.0}{\percent} across the same sweep.

\section{A pilot of direct control against shared autonomy}
\label{sec:pilot-results}

\textbf{Five operators (P1--P5)} completed twenty-three recorded sessions
under the protocol of \cref{sec:ethics}: direct VR control and
shared-autonomy-assisted control, on one of three tasks for P1--P4, and on
an unlabelled task for P5, whose seven sessions predate the per-task naming
convention and are identified as direct/shared from the session's own
logged flag. Coverage is uneven by design -- not every operator performed
every task under both conditions in the time available -- so the pooled
comparison below mixes tasks together and the paired comparison that
follows it is the part of the data that isolates the
condition cleanly. Full session-by-session data, the coverage table, and
the supplementary measures are in \cref{app:pilot}.

\begin{table}[htbp]
  \centering
  \small
  \caption[Pilot session summary]{Twenty-three pilot sessions, pooled by
  condition. Duration includes any pausing or hesitation and is a
  session-length rather than a pure motor-execution measure. Every e-stop
  hold in this pilot occurred once, near the end of the session, to stop
  the recording, and is excluded as a procedural marker rather than a
  safety event.}
  \label{tab:pilot-summary}
  \begin{tabular}{@{}lccc@{}}
    \toprule
    Measure & Direct ($n=13$) & Shared autonomy ($n=10$) \\
    \midrule
    Time to finish, mean (s) & 241 & 219 \\
    Hand travel distance, mean (m) & 2.30 & 2.78 \\
    Re-grips per session, mean & 25.8 & 23.3 \\
    Share of session actively steering, mean & \SI{55}{\percent} & \SI{61}{\percent} \\
    Robot's lag behind the operator's hand, mean (mm) & 2.6 & 2.8 \\
    \bottomrule
  \end{tabular}
\end{table}

\begin{figure}[htbp]
  \centering
  \begin{subfigure}[b]{0.47\linewidth}
    \centering
    \includegraphics[width=\linewidth]{pilot/pilot_time.pdf}
    \caption{pooled by task (P5's task was unrecorded, \emph{P5*})}
    \label{fig:pilot-time-a}
  \end{subfigure}\hfill
  \begin{subfigure}[b]{0.47\linewidth}
    \centering
    \includegraphics[width=\linewidth]{pilot/pilot_paired_time.pdf}
    \caption{paired within operator, both conditions}
    \label{fig:pilot-time-b}
  \end{subfigure}
  \caption[Pilot: time to finish]{Time to finish, pooled
  (\subref{fig:pilot-time-a}) and paired (\subref{fig:pilot-time-b}, the
  fairest comparison here). Duration includes setup and hesitation, so it
  is the least precise pilot measure; \cref{app:pilot}'s supplementary
  measures come straight from the position trace.}
  \label{fig:pilot-time}
\end{figure}

Pooled across all twenty-three sessions (\cref{tab:pilot-summary},
\cref{fig:pilot-time-a}), shared-autonomy sessions averaged shorter
completion times (\SI{219}{\second} against \SI{241}{\second}) and fewer
re-grips per session (\num{23.3} against \num{25.8}), with a robot lag under
\SI{2.9}{\milli\metre} in both conditions -- but \emph{more} hand-travel
distance (\SI{2.78}{\metre} against \SI{2.30}{\metre}), a direction that
holds in the paired comparison too (\cref{sec:pilot-limits}), not just the
pooled one. \Cref{fig:pilot-time-b} isolates the seven cases in which the
\emph{same} operator repeated a task under both conditions: five of seven
finish faster under shared autonomy, and six of seven travel further doing
it.

\begin{figure}[htbp]
  \centering
  \begin{subfigure}[b]{0.31\linewidth}
    \centering\includegraphics[width=\linewidth]{pilot/track_3d.pdf}
    \caption{best VR session, full 3D path}
  \end{subfigure}\hfill
  \begin{subfigure}[b]{0.31\linewidth}
    \centering\includegraphics[width=\linewidth]{pilot/track_joint.pdf}
    \caption{same session, one joint}
  \end{subfigure}\hfill
  \begin{subfigure}[b]{0.34\linewidth}
    \centering\includegraphics[width=\linewidth]{pilot/master_tracking.pdf}
    \caption{master-arm cohort, every session}
  \end{subfigure}
  \caption[Sim-to-real tracking during the pilot]{Every session polled the
  real arms' own encoders alongside the simulated joint state, a live
  sim-to-real check. (a, b) Best-tracking VR session: operator's hand and
  the simulated robot's achieved path, one joint tracking its command to
  \SI{1.06}{\degree} RMS. (c) Three further operators (M1--M3) from the
  mannequin master: two track within \SI{3}{\degree}; the third's three
  sessions track \SIrange{15.9}{21.2}{\degree}, an order of magnitude
  worse and consistent with \cref{fig:channelhealth}. Full record in
  \cref{app:pilot}.}
  \label{fig:track-main}
\end{figure}

This pilot measures operator performance only, at pilot scale, and its
e-stop and lag measures show no evidence of a safety or responsiveness cost
to shared autonomy over the sessions recorded. It does not include a
separate wearer and therefore cannot speak to the wearer-side hypotheses --
trust and embodiment under autonomy -- that motivated the two-person design
of \cref{app:study}; \cref{sec:pilot-limits} treats what would be needed to
extend it into that full comparison.
